\documentclass{article}
\usepackage[margin=1.5cm]{geometry}
\usepackage{amsmath}
\usepackage{enumitem}

\setlist{nosep,leftmargin=*}

\begin{document}
\twocolumn

\title{Problem Set 4: Functions and Modules}
\author{Prof. Jordan C. Hanson}
\date{}

\maketitle

\section{Exercise 1: Comprehensions and \\ Nested Dictionaries}

Suppose we want to search the complex plane over the integer
coordinates

\[
-3 \leq a \leq 3,
\qquad
-3 \leq b \leq 3,
\]

where

\[
z=a+bi.
\]

\begin{enumerate}

\item Use a \textbf{list comprehension} to construct all possible
pairs \([a,b]\):

\begin{verbatim}
points = [
    [a, b]
    for a in range(-3, 4)
    for b in range(-3, 4)
]
\end{verbatim}

How many points should appear in the list?

\vspace{0.3cm}

\item Use a \textbf{dictionary comprehension} to assign an
integer ID to every point:

\begin{verbatim}
complex_grid = {
    i: {
        "real": points[i][0],
        "imag": points[i][1]
    }
    for i in range(len(points))
}
\end{verbatim}

Each value in \texttt{complex\_grid} is itself a dictionary.

Explain why \texttt{complex\_grid} is therefore called a
\textbf{nested dictionary}.

\vspace{0.4cm}

\item Write Python statements that print the real and imaginary
parts associated with key \texttt{10}.

\item Use a loop to print every entry in the form

\begin{verbatim}
ID: 10  z = -2 + 0i
\end{verbatim}

Your exact numerical value for key \texttt{10} will depend on
the order generated by the comprehension.

\end{enumerate}


\section{Exercise 2: Complete the \\ \texttt{ComplexNumber} Class}

Create a Python file named

\begin{verbatim}
complex_number.py
\end{verbatim}

Define the class

\begin{verbatim}
class ComplexNumber:
\end{verbatim}

Each instance must store

\begin{verbatim}
self.real
self.imag
\end{verbatim}

and use an initializer of the form

\begin{verbatim}
def __init__(self, real, imag):
    ...
\end{verbatim}

Implement the following methods.

\subsection*{Addition}

For

\[
z_1=a+bi,
\qquad
z_2=c+di,
\]

\[
z_1+z_2=(a+c)+(b+d)i.
\]

Implement

\begin{verbatim}
def __add__(self, other):
    ...
\end{verbatim}

so that

\begin{verbatim}
z3 = z1 + z2
\end{verbatim}

returns a new \texttt{ComplexNumber} object.

\subsection*{Subtraction}

Use

\[
z_1-z_2=(a-c)+(b-d)i
\]

to implement

\begin{verbatim}
def __sub__(self, other):
    ...
\end{verbatim}

\subsection*{Multiplication}

Complex multiplication is

\[
(a+bi)(c+di)
=
(ac-bd)+(ad+bc)i.
\]

Implement

\begin{verbatim}
def __mul__(self, other):
    ...
\end{verbatim}

\subsection*{Complex Conjugate}

The complex conjugate of

\[
z=a+bi
\]

is

\[
z^*=a-bi.
\]

Implement

\begin{verbatim}
def conjugate(self):
    ...
\end{verbatim}

The method must return a new \texttt{ComplexNumber}.

\subsection*{Division}

Using the conjugate,

\[
\frac{a+bi}{c+di}
=
\frac{ac+bd}{c^2+d^2}
+
\frac{bc-ad}{c^2+d^2}i.
\]

Implement

\begin{verbatim}
def __truediv__(self, other):
    ...
\end{verbatim}

If the denominator is zero, return \texttt{None}.

\subsection*{Display}

Finally, implement

\begin{verbatim}
def __str__(self):
    ...
\end{verbatim}

so that

\begin{verbatim}
z = ComplexNumber(3, -2)
print(z)
\end{verbatim}

produces a readable representation such as

\begin{verbatim}
3 - 2i
\end{verbatim}

Your completed class must contain a docstring describing its
purpose.


\section*{Exercise 3: Finding Poles with \\ the Class}

A rational function

\[
f(z)=\frac{N(z)}{D(z)}
\]

has a \textbf{pole} wherever its denominator is zero:

\[
D(z)=0.
\]

Create a second Python file named

\begin{verbatim}
find_poles.py
\end{verbatim}

Import your class using

\begin{verbatim}
from complex_number import ComplexNumber
\end{verbatim}

Consider

\[
f(z)
=
\frac{1}
{z^2+2z+5}.
\]

We will search the integer complex grid from Exercise 1 to find
values of \(z\) for which

\[
D(z)=z^2+2z+5=0.
\]

\subsection*{Part A: Build the Candidates}

Use a list comprehension to create \texttt{ComplexNumber}
objects:

\begin{verbatim}
candidates = [
    ComplexNumber(a, b)
    for a in range(-3, 4)
    for b in range(-3, 4)
]
\end{verbatim}

\subsection*{Part B: Define the Denominator}

Write

\begin{verbatim}
def denominator(z):
\end{verbatim}

that computes

\[
z^2+2z+5
\]

using only \texttt{ComplexNumber} objects and the overloaded
operators from your class.

For example, define

\begin{verbatim}
two = ComplexNumber(2, 0)
five = ComplexNumber(5, 0)
\end{verbatim}

and construct the expression from these objects.

\subsection*{Part C: Search for Poles}

Loop through \texttt{candidates}. For each candidate,

\begin{verbatim}
value = denominator(z)
\end{verbatim}

A pole has been found when

\begin{verbatim}
value.real == 0 and \
value.imag == 0
\end{verbatim}

Store every pole that you find in a list named

\begin{verbatim}
poles
\end{verbatim}

and print the results.

Verify mathematically that your program should find

\[
\boxed{z=-1+2i}
\]

and

\[
\boxed{z=-1-2i}.
\]

\subsection*{Part D: Store the Results}

Create a dictionary using a dictionary comprehension in which

\begin{itemize}
\item the key is the pole number \(1,2,\ldots\), and
\item the value is a nested dictionary containing the real and
imaginary parts.
\end{itemize}

The structure should have the form

\begin{verbatim}
{
 1: {"real": -1, "imag":  2},
 2: {"real": -1, "imag": -2}
}
\end{verbatim}

Do not type the pole values directly into the dictionary.
Construct the dictionary from the list \texttt{poles} produced
by your search.

\section*{Files to Submit}

Submit exactly the following three files:

\begin{verbatim}
complex_number.py
find_poles.py
problem_set_6.pdf
\end{verbatim}

The PDF should contain your written answers, mathematical work,
and a short explanation of how your pole-search algorithm works.

\end{document}